\documentclass{exam}


%%% Declare commands whether new or old
\newcommand{\declarecommand}[1]{\providecommand{#1}{}\renewcommand{#1}}

%------------------------------------------------------------------------
%	COURSE INFORMATION
%------------------------------------------------------------------------

%%% Global document information
\declarecommand{\university}{University of Toronto}
\declarecommand{\coursename}{Multivariable Calculus with Proofs}
\declarecommand{\coursecode}{MAT237}
\declarecommand{\term}{2026-27}
\declarecommand{\firstterm}{Fall 2026}
\declarecommand{\secondterm}{Winter 2027}

\declarecommand{\title}{\docname}
\declarecommand{\subtitle}{\subname}
\declarecommand{\author}{A. Qiu, S. Vasudevan, and A. Zaman}


%%% URLs
\declarecommand{\urlFAQ}{https://q.utoronto.ca/courses/455325/pages/policies-and-faq} % Quercus FAQ page
\declarecommand{\urlGrading}{https://q.utoronto.ca/courses/455325/external_tools/26926} % Gradescope
\declarecommand{\urlLMS}{https://q.utoronto.ca/courses/455325} % Quercus 


%%%%%%%%%%%%%%%%%%%%%%%%%%%%%%%%%%%%%%%%%%%%%%%%%%%%%%%%%%%%%%%%%%%%%%%%%%%%%%%%%%%%%%%%%%%%%%%%%%%%
%%%		PACKAGES
%%%%%%%%%%%%%%%%%%%%%%%%%%%%%%%%%%%%%%%%%%%%%%%%%%%%%%%%%%%%%%%%%%%%%%%%%%%%%%%%%%%%%%%%%%%%%%%%%%%%

%------------------------------------------------------------------------
%	FORMAT AND LAYOUT
%------------------------------------------------------------------------
\usepackage[24hr,iso]{datetime} % Required to write time of compilation
\usepackage[margin=1in]{geometry} % Required for setting margins
\parindent0pt % Set indentation to 0 pt
\usepackage[english]{babel} % English language/hyphenation
% exam class conflicts with fancyhdr package 

\declarecommand{\toggledraft}{
	\usepackage[firstpageonly=true]{draftwatermark}
	\def\answerborderwidth{1pt}
}

%------------------------------------------------------------------------
%	COLORS
%------------------------------------------------------------------------
\usepackage{xcolor} % required to defined colours
	% Bright qualitative color scheme 
	% Color-blind safe
	% Retrieved 2023-08-27 from Paul Tol (https://personal.sron.nl/~pault/)
	\definecolor{Blue}{HTML}{4477AA}
	\definecolor{Cyan}{HTML}{66CCEE}
	\definecolor{Green}{HTML}{228833}
	\definecolor{Yellow}{HTML}{CCBB44}
	\definecolor{Red}{HTML}{EE6677}
	\definecolor{Purple}{HTML}{AA3377}
	\definecolor{Gray}{HTML}{BBBBBB}
	
	\colorlet{ColorBase}{Purple}
%------------------------------------------------------------------------
%	FONTS
%------------------------------------------------------------------------
\usepackage[charter,cal=cmcal]{mathdesign} % Specify font for serif family
\usepackage{avant} % Specify avant font for sans serif family
\usepackage{mathtools} % Replaces amsmath
\usepackage{sectsty} % Specify section formatting
	\sectionfont{\sffamily\color{ColorBase}}
	\subsectionfont{\sffamily\color{ColorBase!70}}
	\subsubsectionfont{\small\sffamily}

%------------------------------------------------------------------------
%	TABLES, IMAGES, AND FIGURES
%------------------------------------------------------------------------

\usepackage{array} % Required for tables
\usepackage{graphicx} % Required for including images
	\graphicspath{{img/}} % Specifies the directory where pictures are stored
\usepackage{tikz} % Required for creating plots
	\usetikzlibrary{decorations.pathreplacing}
\usepackage{tikz-3dplot} % Required for 3D plots
\usepackage{tikz-cd} % Required for commutative diagrams
\usepackage{pgfplots}
	\pgfplotsset{compat=1.12}
	\usepgfplotslibrary{colormaps}
	\usepgfplotslibrary{patchplots}
	\usepgfplotslibrary{fillbetween}

%------------------------------------------------------------------------
%	LINKS AND REFERENCES
%------------------------------------------------------------------------

\usepackage[hidelinks, urlcolor=Blue, linkcolor=black, colorlinks=true]{hyperref}   
\usepackage[noabbrev,capitalise,nameinlink]{cleveref}

%------------------------------------------------------------------------
%	BOXES, CHOICES, AND LISTS
%------------------------------------------------------------------------

\usepackage{enumerate}
\usepackage{enumitem} % Customize lists
% \setlist{nolistsep} % Reduce spacing between bullet points and numbered lists

\declarecommand{\thequestion}{\sffamily\bfseries Q{\arabic{question}}}
\declarecommand{\partlabel}{(\arabic{question}\alph{partno})} % parts are labelled (1a)
\declarecommand{\thepartno}{(\arabic{question}\alph{partno})}


\usepackage{environ}
%%% Answer environment makes a framed box of a chosen height
\NewEnviron{answer}[2]{
	\fullwidth{
	\vspace*{-10pt}
		\par\nobreak\vspace{\ht\strutbox}\noindent
		\setlength{\fboxrule}{#2}
		\fcolorbox{Gray}{white}{% 
		\parbox[c][#1][t]{\dimexpr\linewidth-2\fboxsep}{
		\hrule width \hsize height 0pt \color{ColorBase} 
		\BODY
		}%
		}%
	}
}

\CorrectChoiceEmphasis{\color{ColorBase}\sffamily\bfseries}	



% Define remark environment
\NewEnviron{remark}{
	\fullwidth{
	\vspace*{-10pt}
		\color{Blue} \textbf{Remark:}  
		\BODY
	}
}


\declarecommand{\CorrectChoiceBox}{
	\checkboxchar{$\Box$} 
	\checkedchar{$\blacksquare$}
}
\declarecommand{\CorrectChoiceCircle}{
	\checkboxchar{\tikz[baseline=-0.6ex]\draw[black,thick] (0,0) circle (1ex);}
	\checkedchar{\tikz[baseline=-0.6ex]\draw[ColorBase,fill=ColorBase,thick] (0,0) circle (1ex);}
}

%------------------------------------------------------------------------
%	HANDWRITING SPACE AND PROMPTS
%------------------------------------------------------------------------

%%% An empty ruled box of fixed height, for handwriting.
%%% Regions are matched in Gradescope, so do not change a height
%%% without re-matching the outline.
\declarecommand{\ruledbox}[1]{%
	\par\nobreak\noindent
	\begingroup
	\setlength{\fboxrule}{0.5pt}%
	\setlength{\fboxsep}{6pt}%
	\fcolorbox{Gray}{white}{%
		\parbox[c][#1][t]{\dimexpr\textwidth-2\fboxsep-2\fboxrule\relax}{%
			\rule{\linewidth}{0pt}% forces the box to full width; an empty parbox has zero width
		}%
	}%
	\endgroup
	\par
}


%%% A labelled pre-meeting prompt. Keep the box small on purpose:
%%% a full page of prose here costs the TA the 60-second skim.
\declarecommand{\prompt}[2]{%
	\vspace{6pt}
	\par\noindent{\bfseries\color{ColorBase} #1}~{#2}
}

%%% A labelled rule for handwritten fields.
\declarecommand{\field}[2]{{\sffamily\small #1}~\raisebox{2pt}{\makebox[#2]{\hrulefill}}}


%------------------------------------------------------------------------
%	THEOREMS
%------------------------------------------------------------------------

\usepackage{amsthm}

%%% Specify theorems
\newtheorem{theorem}{Theorem}
\newtheorem*{definition}{Definition}

\newtheoremstyle{named}{}{}{\itshape}{}{\bfseries}{.}{.5em}{\thmnote{#3}}
\theoremstyle{named}
\newtheorem*{namedtheorem}{Theorem}


%% Setup proof lines
\NewEnviron{lines}{
	\vspace*{-5pt}
	\begin{center}
	\fbox{\begin{minipage}{0.9\textwidth}
	\begin{enumerate}[label=\arabic*.]
		\it \BODY
	\end{enumerate}
	\end{minipage}
	}
	\end{center}
}
 
%------------------------------------------------------------------------
%	HEADER & FOOTER
%------------------------------------------------------------------------

%%% Define footnote without marker
\declarecommand{\blfootnote}[1]{%
  \begingroup
  \declarecommand\thefootnote{}\footnote{#1}%
  \addtocounter{footnote}{-1}%
  \endgroup
}

%%% Specify header and footer
\pagestyle{headandfoot} 	\runningfootrule  
	\runningfooter{\coursecode}{\docname\ - Page \the\numexpr\thepage\ of \numpages}{\term}

%------------------------------------------------------------------------
%	TITLE PAGE
%------------------------------------------------------------------------


%%%% Define fancy title
\declarecommand{\maketitle}{
\thispagestyle{empty} % Suppress headers and footers on the title page
\vspace*{-85pt}
\hspace*{-90pt}
\begin{tikzpicture}[remember picture]
\draw (8,-2) node [fill=ColorBase!20,fill opacity=0.8,text opacity=1,inner sep=15pt]{\Huge\centering\parbox[c][][t]{\paperwidth}{\centering\huge\color{ColorBase!60}\bfseries\sffamily\coursecode\, \coursename~\\\Huge\color{ColorBase}\title\\ % Title
{\Large\color{black}\bfseries\sffamily\subtitle}}}; % Subtitle
	\end{tikzpicture}
	\vspace*{-10pt}
	\blfootnote{Created by \author. Last updated \today \, at \currenttime.} 
}

%%% Uncomment if you are writing a draft
% \toggledraft

%%% Uncomment if you want to show your multiple choice answers
\printanswers

%%% Local document information
\declarecommand{\docname}{Problem Set 3 Debrief}
\declarecommand{\dueday}{Friday}
\declarecommand{\duedate}{November 20, 2026}
\declarecommand{\duetime}{23:59 ET}
\declarecommand{\subname}{Due \dueday\,\duedate\,at\,\duetime}

%%% Student information
\declarecommand{\FirstStudentName}{Sebastian Cooper-Giannotta}
\declarecommand{\SecondStudentName}{Yueheng Guan}
% sign the PDF after compiling

%%% Custom math commands
%% Standard
\declarecommand{\ds}{\displaystyle} % force math styling
\declarecommand{\emptyset}{\varnothing} % better empty set
\declarecommand{\epsilon}{\varepsilon} % better epsilon
\declarecommand{\R}{\mathbb{R}} % reals
\declarecommand{\Z}{\mathbb{Z}} % integers
\declarecommand{\Q}{\mathbb{Q}} % rationals
\declarecommand{\N}{\mathbb{N}} % naturals
\declarecommand{\C}{\mathbb{C}} % complex
\DeclarePairedDelimiter\abs{\lvert}{\rvert} % absolue value
\DeclarePairedDelimiter\norm{\lVert}{\rVert} % norm

%% Linear algebra
\declarecommand{\mat}[1]{\begin{bmatrix*}[r]#1\end{bmatrix*}}
\declarecommand{\matc}[1]{\begin{bmatrix}#1\end{bmatrix}} % matrix
\DeclareMathOperator{\Span}{span} % span
\DeclareMathOperator{\Img}{img} % image
\DeclareMathOperator{\Ker}{ker} % kernel
\DeclareMathOperator{\Id}{id} % identity map
\DeclareMathOperator{\Range}{range} % range
\DeclareMathOperator{\Rref}{rref} % row reduced echelon form
\DeclareMathOperator{\Rank}{rank} % rank
\DeclareMathOperator{\Null}{null} % null space
\DeclareMathOperator{\Nullity}{nullity} % nullity
\DeclareMathOperator{\Proj}{proj} % projection
\DeclareMathOperator{\Dim}{dim} % dimension

%% Topology
\DeclareMathOperator{\Int}{int}
\DeclareMathOperator{\Cl}{cl}

%% Calculus
\DeclareMathOperator{\Length}{length} % length
\DeclareMathOperator{\Area}{area} % area
\DeclareMathOperator{\Vol}{vol} % volume
\DeclareMathOperator{\Grad}{grad} % gradient
\DeclareMathOperator{\Curl}{curl} % curl
\DeclareMathOperator{\Div}{div} % divergence

\begin{document}

\maketitle

\bigskip

%%%%%%%%%%%%%%%%%%%%%%%%%%%%%%%%%%%%%%%%%%
%%%%%%%%% COVER PAGE
%%%%%%%%%%%%%%%%%%%%%%%%%%%%%%%%%%%%%%%%%%
%%% HEADER BLOCK
{\sffamily We,  the undersigned students, were present during the debrief and contributed to this write-up.}\\
\setlength{\fboxrule}{0.8pt}
\noindent\parbox[c][96pt][c]{\dimexpr\textwidth-2\fboxsep-2\fboxrule}{
	\vspace{2pt}
	\makebox[0.6\textwidth]{\sffamily Student 1: {\color{ColorBase}\bfseries \FirstStudentName} \hrulefill } \,\,
    \makebox[0.4\textwidth]{\sffamily Signature: \hrulefill}

	\vspace{14pt}

    \makebox[0.6\textwidth]{\sffamily Student 2: {\color{ColorBase}\bfseries \SecondStudentName} \hrulefill } \,\,
    \makebox[0.4\textwidth]{\sffamily Signature: \hrulefill}

	\vspace{16pt}

	\makebox[0.95\textwidth]{\sffamily Which problem are you debriefing? Select ONE. \CorrectChoiceCircle
	\hspace{10pt}
    % Replace \chocie with \CorrectChoice for your preferred selection
	\begin{oneparcheckboxes}
		\choice{\bfseries Q1} \qquad 
	    \choice {\bfseries Q2} \qquad
	    \CorrectChoice {\bfseries Q3} \qquad
	\end{oneparcheckboxes} 
		
	}
	\vspace{2pt}
}

\bigskip
%%% HEADER BLOCK
\setlength{\fboxrule}{1pt}



\fbox{
\begin{minipage}{\textwidth}
\colorbox{ColorBase!20}{%
	\parbox[l][30pt][l]{\dimexpr\textwidth-2\fboxsep}{
	\centering \sffamily {\color{ColorBase}\bfseries\Large TA Debrief Feedback} \\ 
	{\bfseries This box is only to be filled out by your TA.} } 
}

\bigskip
\sffamily
\begin{tabular}{ll}
\prompt{1.}{Does the submission contain a detailed attempted solution?} & 
\begin{oneparcheckboxes}
	\choice Yes
	\choice No	
\end{oneparcheckboxes} \\
\prompt{2.}{Did both students participate in the discussion?} & 
\begin{oneparcheckboxes}
	\choice Yes
	\choice No	
\end{oneparcheckboxes} \\
\end{tabular}

\bigskip
\sffamily
\begin{minipage}{0.2\textwidth}
\begin{checkboxes}
	\choice Brendan
	\choice Dai
\end{checkboxes}
\end{minipage}
\begin{minipage}{0.2\textwidth}
\begin{checkboxes}
	\choice Henry
	\choice Joshua
\end{checkboxes}
\end{minipage}
\begin{minipage}{0.2\textwidth}
\begin{checkboxes}
	\choice Nick
	\choice Nikolas
\end{checkboxes}
\end{minipage}
\begin{minipage}{0.2\textwidth}
\begin{checkboxes}
	\choice Ren\'{e}e S
	\choice Renee W
\end{checkboxes}
\end{minipage}
\bigskip 


\bigskip 
\makebox[0.95\textwidth]{TA Signature : \enspace{}\hrulefill}		
	
{\sffamily 
\bigskip 
\hspace*{-10pt}

The feedback below does \textbf{not} affect your grade. The TA will give you the rubric during the debrief.\\ 

\begin{tabular}{ll}
\prompt{3.}{They verbally explained their  solution today.} & 
\begin{oneparcheckboxes}
	\choice Fully
	\choice Partly
	\choice Not this time
\end{oneparcheckboxes} \\
\prompt{4.}{The solution satisfies Rubric Item A (structure).} & 
\begin{oneparcheckboxes}
	\choice Fully
	\choice Partly
	\choice Not this time 
\end{oneparcheckboxes} \\
\prompt{5.}{The solution satisfies Rubric Item B (key idea).} & 
\begin{oneparcheckboxes}
	\choice Fully
	\choice Partly
	\choice Not this time 
\end{oneparcheckboxes} \\
\prompt{6.}{The solution satisfies Rubric Item C (writing).} & 
\begin{oneparcheckboxes}
	\choice Fully
	\choice Partly
	\choice Not this time
\end{oneparcheckboxes} \\
\end{tabular}
}
\end{minipage}
}

\bigskip 

\prompt{IV.}{\textbf{After your debrief}, summarize at least one takeaway from the TA's feedback on your solution.}{}
\begin{answer}{100pt}{1pt} % Do not change the height of this box
    % Insert your answer here. Text should be purple.
\end{answer}

\vfill

\pagebreak 
%%% THREE PROMPTS
\section*{Your discussion points}

\textbf{Fill these in before your debrief.} Your TA may or may not use these points to guide discussion.

\prompt{I.}{What was your first idea that did not work, and how did you know?}{}

\begin{answer}{100pt}{1pt} % Do not change the height of this box
    % Insert your answer here. Text should be purple.
\end{answer}

\bigskip
\prompt{II.}{Pick one theorem you cited. State its hypotheses and point to the line where you checked each.}

\begin{answer}{100pt}{1pt} % Do not change the height of this box
    % Insert your answer here. Text should be purple.
\end{answer}

\bigskip

\prompt{III.}{Select ONE question you could ask the TA.}

\medskip
\CorrectChoiceCircle
% Replace \chocie with \CorrectChoice for your preferred selection
\begin{checkboxes} \sffamily
	\choice Does our argument have all the main ideas?
	\choice Did we state everything precisely, e.g. variables, quantifiers, definitions?
	\choice Does our argument have an easy-to-read logical flow?
	\choice Do we use proper notation and terminology?
	\choice Could we add more justification somewhere?  
\end{checkboxes}

\vfill


\pagebreak
%%%%%%%%%%%%%%%%%%%%%%%%%%%%%%%%%%%%%%%%%%
%%%%%%%%% SOLUTION
%%%%%%%%%%%%%%%%%%%%%%%%%%%%%%%%%%%%%%%%%%
\section*{Your solution} 
{
Write a polished solution to your team's chosen question from page 1, including all parts. Mistakes welcome.
\smallskip
\vfill
\begin{answer}{600pt}{0pt} % Do not change the height of this box
    % Insert your answer here on page 3. Text should be purple.
    \textbf{(3a)}
    \setlength{\parskip}{1em}
    The idea is to show that F is equivalent to a C1 function.

    first, take $\pi_1 : \mathbb{R}^k \times \mathbb{R}^{n-k} \to \mathbb{R}^k$ by $(x,y) \to x$ similarly for projection on $\pi_2$.

    Both $\pi_1$ and $\pi_2$ are linear and thus are $C^1$.

    Then, noting that $F$ is $C^1$, by Corollary 4.1.5, we have that $F\circ \pi_1 : \mathbb{R}^k \times \mathbb{R}^{n-k} \to \mathbb{R}^k \to \mathbb{R}^{n-k}$ is also $C^1$.

    And, since $\pi_2$ is $C^1$, apply theorem 3.6.6 and we get that $\pi_2- F \circ \pi_1$ is $C^1$. Then by definition, G which has the same domain/codomain and maps (x,y) to the same values as this function, is $C^1$.
\end{answer}


\pagebreak
{\centering\footnotesize (Extra space provided for your solution. Leave blank if you don't need it.)}
\begin{answer}{600pt}{0pt} % Do not change the height of this box
    % Continue your answer here on page 4. Text should be purple.

    \textbf{(3b)}
    \setlength{\parskip}{1em}
    $DG(p) = [\partial_1G(p),...,\partial_nG(p)] = [\frac{\partial G}{\partial x_1}(p),...,\frac{\partial G}{\partial x_k}(p),\frac{\partial G}{\partial y_1}(p),...,\frac{\partial G}{\partial y_{n-k}}(p)]$
    
    But since y is a constant w.r.t x1 through xk, we know that the first k partial derivatives are only w.r.t F, and the last n-k are y w.r.t its own components:
    
    $=[\frac{-\partial F}{\partial x_1}(a),...,\frac{-\partial F}{\partial x_k}(a),e_1,e_2,...,e_{n-k}]$
    since F takes input x in domain
    
\end{answer}
~

\vfill
\pagebreak
{\centering\footnotesize (Extra space provided for your solution. Leave blank if you don't need it.)}
\begin{answer}{600pt}{0pt} % Do not change the height of this box
    \textbf{(3c)}

    \begin{proof}
        \setlength{\parskip}{1em}

        Let $S = \{(x, y) \in \R^k \times \R^{n - k} : y = F(x)\}$.

        Since $S$ is the graph of $F$, and $F$ is $C^1$, we have by the tangent space of graphs theorem:
        \[
            T_pS = \{(h, dF_a(h)) : h \in \R^k\}
        \]

        Let $(x, y) \in \R^k \times \R^{n - k}$, then:
        \begin{align*}
            (x, y) \in \Ker(dG_p) & \iff (x, y) \in \Null(DG(p)) \\
            & \iff DG(p) \mat{x \\ y} = \mat{-DF(a) & I_{n - k}} \mat{x \\ y} = 0 \\
            & \iff -DF(a)x + I_{n - k}y = y - DF(a)x = 0 \\
            & \iff y = DF(a)x = dF_a(x) \\
            & \iff (x, y) = (x, dF_a(x)) \\
            & \iff (x, y) \in T_pS
        \end{align*}

        Therefore $T_pS = \Ker(dG_p)$.
    \end{proof}
\end{answer}
~

\end{document}
